\section{Lower reductions: certified quartics carry the full sign problem}
\label{sec:reductions}

Theorem~\ref{thm:exact-upper} turns every exact order or equality test at
the minimizer of an explicitly given, globally strongly convex polynomial
into one $\PosSLP$ instance. This section shows that the order tests are no
easier, even in a setting that removes every familiar source of difficulty:
a single quartic on all of $\R^n$, supplied with a rational positive definite
Hessian Gram matrix that certifies uniform strong convexity and can be checked
in polynomial time. The tests stay $\PosSLP$-hard when the optimal value is
promised to be nonzero. The coordinate tests stay $\PosSLP$-hard when the
optimizer is promised to be rational, to lie in $[-1,1]^n$, and to have the
known optimal value zero.

These are the exact decisions that a branch-and-bound method for
mixed-integer nonlinear optimization makes at a node: whether the bound of a
node relaxation lies below an incumbent value, and on which side of a
branching threshold an optimal coordinate lies. Certified strong convexity,
the absence of constraints, and a rational optimizer each remove a known
obstacle to exact computation, but none of them removes the arithmetic one.
An ordinary polynomial-time algorithm for any of the comparisons below would
place $\PosSLP$ in P. \citet{AllenderEtAl2009} show that $\PosSLP$ lies in the
counting hierarchy and that $\mathrm{P}^{\PosSLP}$ is the Boolean part of
polynomial-time computation over the reals with constants $0$ and $1$ in the
Blum--Shub--Smale model; whether $\PosSLP$ lies in P is open. The results here
concern exact comparisons only. Approximate optimal values and approximate
optimizers of the same quartics are computable in polynomial time
(Theorem~\ref{thm:global-point}).

\paragraph{Mechanism.}
Every lower bound below has three steps.
\begin{enumerate}[label=(\arabic*)]
\item A \emph{sign compiler} converts an integer straight-line program into a
circuit whose gate values stay bounded and whose designated output has the
sign of the program's output. The compilers never represent the integer
values themselves, which can have doubly exponential magnitude. They carry
only signs, through real signals that can be as small as $2^{-2^{cT}}$ for a
program with $T$ operations.
\item A \emph{realization} places the list of all gate values at the unique
zero of one rational quartic. A gate-specific exposing quadratic supplies
positive curvature in every direction at that zero, and the assembly
Lemma~\ref{lem:quartic-realization} combines it with the gate residuals into
a globally strongly convex sum of squares with a rational positive definite
Hessian Gram.
\item For value comparisons, a \emph{cubic perturbation} transfers the
designated sign to the sign of the optimal value without losing the Hessian
certificate.
\end{enumerate}
We use two compilers. Certified cube-root circuits
(Section~\ref{sec:reductions-root-compiler}) give hardness of exact
feasibility, of optimization over a fixed cube, and of the sign of an
unconstrained minimum. Circuits of rational unit quaternions
(Section~\ref{sec:reductions-quaternion}) keep every gate value rational and
give coordinate hardness with a rational optimizer.

\subsection{Statements}
\label{sec:reductions-statements}

$\PosSLP$ asks whether the output $V$ of a division-free integer
straight-line program with constants $0,1$ and operations $+,-,\times$ is
positive. Programs and all circuits below are directed acyclic graphs: a gate
may be used by any number of later gates, and sizes count gates of the shared
graph, never an expanded formula or expanded fractions. The integer $W=2V-1$
is never zero, and $W>0$ exactly when $V>0$; replacing $V$ by $1-V$ reverses
the answer. All reductions are deterministic polynomial-time many-one
reductions, and $L$ denotes total binary input length.

A \emph{certified quartic} on $\R^n$ is a pair $(f,M)$, where
$f\in\Q[x_1,\ldots,x_n]$ has degree at most four and $M$ is a rational
symmetric positive definite matrix of order $n+n^2$ such that
\begin{equation}
 v^{\mathsf T}\nabla^2f(x)\,v=(v,x\otimes v)^{\mathsf T}M\,(v,x\otimes v)
 \qquad\text{identically in }(x,v)\in\R^n\times\R^n,
 \label{eq:reductions-certified}
\end{equation}
with $(x\otimes v)_{(k,j)}=x_kv_j$ (Definition~\ref{def:models-hessian-gram}).
Coefficient comparison and rational symmetric elimination check this format in
polynomial time. The matrix certifies $\nabla^2f\succeq\lambda_{\min}(M)I$, and
$\det M/(\tr M)^{n+n^2-1}$ is a positive rational lower bound on
$\lambda_{\min}(M)$ of polynomial bit length
(Lemma~\ref{lem:models-det-trace}). Hence $f$ has a unique minimizer $p$ and
optimal value $f^*=f(p)$.

\begin{theorem}[Certified quartic classification]
\label{thm:quartic-complete}
Consider inputs consisting of a certified quartic $(f,M)$ on $\R^n$, a
rational threshold $r$ and, for coordinate tests, an index
$j\in\{1,\ldots,n\}$. Inputs violating the format are no-instances.
\begin{enumerate}[label=(\roman*)]
\item Each of the eight tests $f^*<r$, $f^*\le r$, $f^*>r$, $f^*\ge r$,
$p_j<r$, $p_j\le r$, $p_j>r$, $p_j\ge r$ is $\PosSLP$-complete.
\item Each of the tests $f^*=r$ and $p_j=r$ reduces to $\PosSLP$. No matching
lower bound is asserted.
\item Each of the following properties is $\PosSLP$-complete: $f\ge0$ on
$\R^n$; $f$ is a sum of squares of real polynomials; $f$ has a rational
positive definite Gram matrix on the vector of all monomials of degree at most
two. Being a sum of squares of rational polynomials is $\PosSLP$-hard.
\item The lower bounds for the value tests in (i) and for all properties in
(iii) hold already for instances with $r=0$, $M\succeq I$ and $f^*\ne0$. The
lower bounds for the coordinate tests hold under the promises of
Theorem~\ref{thm:rational-optimizer}.
\end{enumerate}
\end{theorem}

The upper bounds are Theorem~\ref{thm:exact-upper}; the content of this
section is the lower bounds. Part~(iii) concerns three different certificate
notions. In the certified class, nonnegativity
and real sums of squares are equivalent, and a rational positive definite
polynomial Gram exists exactly when $f^*>0$. Rational sums of squares at
$f^*=0$ are a different matter: Example~\ref{ex:fields-ternary} is a
certified quartic with minimum zero that is a sum of squares over a real field
$E$ exactly when $2^{1/5}\in E$. No upper bound for rational sum-of-squares
membership is claimed.

Quartics are the lowest degree at which these lower bounds can hold. A
globally convex polynomial of degree at most three is quadratic
(Lemma~\ref{lem:models-low-degree}), and a strongly convex rational quadratic
has a rational minimizer computable by rational linear algebra.

\begin{theorem}[Rational optimizer and known value]
\label{thm:rational-optimizer}
There is a polynomial-time many-one reduction from $\PosSLP$ to the test
$p_j>0$ for certified quartics, whose output instances $(F,M,j)$ on $\R^N$
satisfy all of the following promises:
\begin{enumerate}[label=(\roman*)]
\item $p\in\Q^N\cap[-1,1]^N$ and $p_j\ne0$;
\item $F^*=0$, and $F$ is supplied as a sum of $N+1$ squares of rational
quadratic polynomials;
\item $\nabla^2F\succeq\frac32I$ on $\R^N$.
\end{enumerate}
Consequently each of the four coordinate order tests at threshold zero remains
$\PosSLP$-complete on the instances satisfying these promises.
\end{theorem}

Rationality of the optimizer is a promise satisfied by the constructed
instances; the reduction neither recognizes nor needs to recognize it. The
theorem shows that exact field membership does not supply an exact algorithm:
the optimizer is rational, its value is known, and the objective comes with
short rational square factors, yet a single coordinate sign still carries an
arbitrary $\PosSLP$ instance. Rationality also gives no bound on representation
size: in the constructed instances one coordinate of $p$ is a nonzero rational of absolute
value below $2^{-2^{cN}}$ for an absolute constant $c>0$
(Remark~\ref{rem:reductions-rational-length}), so its expanded fraction has
exponentially many bits. Theorem~\ref{thm:rational-height} gives
unconditional families of bounded rational optimizers with long expanded
encodings and additional conditioning properties. Corollary~\ref{cor:constraints-binary}
uses the instances of Theorem~\ref{thm:rational-optimizer} to show that one
optimal binary decision remains $\PosSLP$-hard to extract when the optimal
value is known in advance.

\paragraph{Relation to earlier hardness results.}
Exact convex feasibility was already known to capture arithmetic circuit
comparison: \citet{TarasovVyalyi2008} reduce $\PosSLP$ to exact feasibility of
semidefinite programs. Bounded arithmetic circuits also capture $\PosSLP$:
\citet[Lemma~5]{EY2010} reduce it, with a circuit of linear size over
$\{+,\times,/\}$ whose gate values all lie in $(0,1)$, to an order comparison
of that circuit's output.
The contribution here is the restriction to one explicitly expanded
quartic with a checked strict Hessian certificate and no constraints, together
with the two sign compilers it requires. The analytic compiler avoids products
of independent values, which the quartic realization cannot accept; the
quaternion compiler keeps every gate value exactly rational and bounded.
Commutators of near-identity rotations are a classical tool for multiplying
small rotation angles, used in the Solovay--Kitaev algorithm for approximate
synthesis in $\mathrm{SU}(2)$~\cite{DawsonNielsen2006}, and identity testing
for shared circuits over linear groups has also been
studied~\cite{KoenigLohrey2015}.
Those works concern approximate synthesis and identity testing,
respectively; they do not decide the sign of a designated rational coordinate
in a fixed compact group. Our use requires exact rational gates and error
control below a doubly exponentially small nonzero signal.
Version~1 of \citet{SlotSteurerWiedmer2025} proves polynomial-time
approximation for convex polynomial programming and leaves the complexity of
the exact decision problem $\exists x\in P:\ f(x)\le0$ for convex quartics
open.
Corollary~\ref{cor:reductions-feasibility} shows that this decision is
$\PosSLP$-hard already when $P$ is a halfspace or the unit cube and $f$ is a
certified quartic. It does not determine the complexity class of that
constrained problem: Theorem~\ref{thm:exact-upper} concerns unconstrained
minimization, and Section~\ref{sec:constraints} treats constrained exact
comparison.

\subsection{Two obstacles}
\label{sec:reductions-obstacles}

The following two facts explain the shape of the constructions. The first
shows why the compilers carry only signs through bounded signals. The second
shows why the realizations need a gate-specific exposing quadratic rather than
squared gate equations.

\begin{proposition}[Radius and naive lifting]
\label{prop:reductions-obstacles}
\leavevmode
\begin{enumerate}[label=(\roman*)]
\item If $F\colon\R^n\to\R$ is twice continuously differentiable with
$\nabla^2F\succeq I$ and minimizer $p$, then $\norm{p}\le\norm{\nabla F(0)}$.
In particular, a rational polynomial whose coefficients have bit length at
most $b$ has a minimizer of norm at most $\sqrt n\,2^b$.
\item For every $\varepsilon>0$ the nonnegative quartic
$H(x,y)=x^2+x^4+\varepsilon(y-x^2)^2$ has the unique zero $(0,0)$, at which
the residual map $(x,y-x^2)$ has the identity Jacobian, but
$\partial_x^2H(0,y)=2-4\varepsilon y$ is negative for $y>1/(2\varepsilon)$.
\end{enumerate}
\end{proposition}

\begin{proof}
For~(i), strong monotonicity of $\nabla F$ gives
$(\nabla F(0)-\nabla F(p))^{\mathsf T}(0-p)\ge\norm{p}^2$; since
$\nabla F(p)=0$, the Cauchy--Schwarz inequality gives the bound. The linear
coefficients of $F$ form $\nabla F(0)$. For~(ii), the zero set is clear, and
$\partial_x^2[\varepsilon(y-x^2)^2]=-4\varepsilon y+12\varepsilon x^2$.
\end{proof}

Part~(i) rules out realizations that keep unscaled integer values: a
program with $m$ squarings computes $2^{2^m}$, whereas a polynomial-size
normalized quartic has a minimizer of singly exponential norm. Part~(ii) shows
that a nonsingular Jacobian at the intended zero does not make squared gate
residuals globally convex: at $x=0$ the squared residual term
$\varepsilon(y-x^2)^2$ contributes $-4\varepsilon y$ to $\partial_x^2H$, which
is unbounded below as $y$ grows. The realizations below control such negative
curvature with one additional square, that of a gate-specific exposing
quadratic whose quadratic part is positive definite; after scaling, this
square makes the whole sum uniformly strongly convex on the whole space
(Lemma~\ref{lem:quartic-realization}).
This is a feature of our construction, not a necessary condition for
convexity: $(x^2-y^2)^2+(2xy)^2=(x^2+y^2)^2$ is convex, although both squared
quadratics have indefinite quadratic parts.

\subsection{Sign compiler I: certified cube-root circuits}
\label{sec:reductions-root-compiler}

\begin{definition}[Certified odd-root circuits]
\label{def:reductions-root-circuit}
A \emph{certified odd-root circuit} with $k$ gates consists of odd integers
$d_i=2n_i-1\ge3$ written in unary, rational coefficients $c_i$ and $a_{ije}$
for $1\le j<i\le k$ and $1\le e\le n_j$, and rational boxes
$[\underline\xi_i,\overline\xi_i]$ not containing zero, such that for every
$i$
\begin{equation}
 c_i+\sum_{j<i}\sum_{e=1}^{n_j}a_{ije}R_{j,e}
 \subseteq[\underline\xi_i^{\,d_i},\overline\xi_i^{\,d_i}],
 \qquad
 R_{j,e}=\{s^e:\ s\in[\underline\xi_j,\overline\xi_j]\},
 \label{eq:reductions-interval-test}
\end{equation}
where the left side is a Minkowski sum of intervals. The circuit defines reals
$\xi_1,\ldots,\xi_k$ by
\begin{equation}
 \xi_i^{d_i}=\rho_i(\xi):=c_i+\sum_{j<i}\sum_{e=1}^{n_j}a_{ije}\xi_j^e,
 \label{eq:reductions-root-gate}
\end{equation}
taking the unique real $d_i$-th root. It is a \emph{cube-root circuit} if every
$d_i=3$; then each radicand is a rational affine combination of earlier roots
and their squares.
\end{definition}

The test~\eqref{eq:reductions-interval-test}, which we call the
\emph{direct interval test}, uses exact rational arithmetic and unary
exponents, so it is checkable in polynomial time. By induction over $i$, it
implies $\xi_i\in[\underline\xi_i,\overline\xi_i]$ for every gate: if all
earlier roots lie in their boxes, then $\rho_i(\xi)$ lies in the left side of
\eqref{eq:reductions-interval-test}, and the real $d_i$-th root is increasing.
Radicand coefficients may have either sign. A product $\xi_j\xi_l$ of two
different earlier gates is not available; this restriction is what makes a
quartic realization possible, and it is the reason a new compiler is needed.

\begin{theorem}[Cube-root sign compiler]
\label{thm:reductions-root-circuit}
Given an integer straight-line program with output $V$, one can compute in
polynomial time a certified cube-root circuit with boxes
$[1-w_i,1+w_i]$, $0<w_i\le10^{-9}$, and two distinct gates $o$ and $e$ such
that
\begin{equation}
 \xi_o\ne1,\qquad \operatorname{sign}(\xi_o-1)=\operatorname{sign}(2V-1),
 \qquad\abs{\xi_o-1}\le2^{-29},
 \qquad 0<(\xi_e-1)^2\le\tfrac18\abs{\xi_o-1}.
 \label{eq:reductions-compiler-output}
\end{equation}
Every radicand is a rational constant, an affine function
$1+\sum_j\lambda_j(\xi_j-1)$ with $\sum_j\abs{\lambda_j}\le6$, or the square
gate $3\xi_j^2-6\xi_j+4=1+3(\xi_j-1)^2$.
\end{theorem}

The proof is in Appendix~\ref{app:reductions-root-compiler}. Its gadgets act
on small real signals $x$, stored as roots $\xi=1+x$. Put
\[
 \phi(z)=(1+3z)^{1/3}-1,\qquad
 \operatorname{sq}(z)=\phi(z^2),\qquad
 \theta(z)=\tfrac12\bigl[\operatorname{sq}(\phi(z))+\operatorname{sq}(\phi(-z))\bigr],
\]
\[
 \operatorname{mul}(x,y)=\phi\Bigl(\tfrac14\bigl[\theta(x+y)-\theta(x-y)\bigr]\Bigr).
\]
One root gate computes $\phi$ of an affine combination of signals, and one
square gate computes $\operatorname{sq}$; $\operatorname{mul}$ uses nine root
gates. The function $\theta$ is even with $\theta(z)=z^2+O(z^4)$, so
$\operatorname{mul}(x,y)=xy+O(\abs{(x,y)}^4)$. The decisive property is exact:
$\operatorname{mul}(x,0)=\operatorname{mul}(0,y)=0$. A Cauchy estimate then gives
the relative error bound
$\abs{\operatorname{mul}(x,y)-xy}\le2^{24}\abs{xy}(\abs x+\abs y)$, which stays
proportional to the product even when the two inputs have very different
magnitudes. Each integer gate value $V_i$ is carried by a pair of signals of a
common order $d_i$ in a tiny parameter $\delta$, with leading coefficients
$V_i$ and $1$; addition cross-multiplies by the denominator signals, so that
it only ever adds signals of equal order. All error bounds are absolute
multiples of $\delta^{d_i+1}$ and never divide by a leading coefficient, so
cancellation and zero intermediate values need no special treatment. The
parameter $\delta$ is the output of $2T+5$ chained square gates starting from
the printed rational $1000^{-O(T)}$; it is never printed itself. The boxes are
verified directly, independently of the analytic estimates, by
Lemma~\ref{lem:reductions-small-signal}.

For completeness of the picture, the language produced by the compiler is
itself $\PosSLP$-complete.

\begin{theorem}[Certified cube-root comparison]
\label{thm:reductions-root-language}
For a certified cube-root circuit with positive boxes, a gate $o$, and a
rational $r$, each of the tests $\xi_o>r$, $\xi_o\ge r$, $\xi_o<r$, $\xi_o\le r$
is $\PosSLP$-complete, and $\xi_o=r$ reduces to $\PosSLP$. Malformed inputs,
including inputs failing the direct interval test, are no-instances.
\end{theorem}

The lower bounds follow from Theorem~\ref{thm:reductions-root-circuit} with
$r=1$, after replacing $V$ by $1-V$ for the tests $<$ and $\le$. The upper
bound, proved in Appendix~\ref{app:reductions-root-language}, uses the same
architecture as Theorem~\ref{thm:exact-upper} with a more elementary separation
bound: after clearing denominators every gate value is an algebraic integer of
degree at most $3^k$, so a field norm bounds $\abs{\xi_o-r}$ away from zero,
and polynomially many exact Newton steps for each cube root reach that
accuracy in a shared rational circuit.

\subsection{Quartic realization of certified root circuits}
\label{sec:reductions-realization}

\begin{theorem}[Signed odd-root realization]
\label{thm:reductions-signed-root}
There is a deterministic polynomial-time algorithm that, given a certified
odd-root circuit, computes signs $\sigma_i\in\{-1,1\}$, a rational
$\kappa\ge1$, and a rational quartic $F$ in the $N=\sum_in_i$ variables
$X=(X_{i,e})_{1\le i\le k,\,1\le e\le n_i}$ with the following properties.
\begin{enumerate}[label=(\roman*)]
\item $F$ is supplied as a sum of $N+1$ squares of rational quadratic
polynomials, and $\nabla^2F\succeq\frac32I$ on $\R^N$.
\item $F^{-1}(0)=\{p\}$, where $p_{i,e}=(\sigma_i\kappa\xi_i)^e$. In
particular $\xi_i=\sigma_ip_{i,1}/\kappa$.
\item A rational positive definite Hessian Gram matrix for $F$ in the sense
of~\eqref{eq:reductions-certified} is computed as part of the output.
\end{enumerate}
Here $\sigma_i$ is the sign of the box of gate $i$, and
$\kappa=\max\{1,\max_i1/\min(\abs{\underline\xi_i},\abs{\overline\xi_i})\}$.
\end{theorem}

The proof is in Appendix~\ref{app:reductions-signed-root}. After the
normalization $\alpha_i=\sigma_i\kappa\xi_i\ge1$, the coordinates $X_{i,e}$
stand for the powers $\alpha_i^e$, $e\le n_i$, and the $N$ residuals
\[
 X_{i,1}X_{i,e}-X_{i,e+1}\quad(1\le e<n_i),\qquad
 X_{i,n_i-1}X_{i,n_i}-b_i(X),\quad
 b_i(X)=c_i'+\sum_{j<i,\,e}a_{ije}'X_{j,e},
\]
have $p$ as their unique real zero, with a block-triangular Jacobian whose
determinant has absolute value $\prod_id_i\alpha_i^{d_i-1}\ge1$. For one gate with root $\alpha$ and
$d=2n-1$, put $\ell_j=X_{j+1}-\alpha X_j$ for $0\le j<n$, with $X_0=1$. A
tridiagonal positive definite form $P_\alpha=\ell^{\mathsf T}T_\alpha\ell$
restricts to the power curve $X_j=t^j$ as
\[
 P_\alpha(t,t^2,\ldots,t^n)=(t-\alpha)(t^d-\alpha^d),
\]
and the gate's exposing quadratic is
\[
 E_i^*=P_{\alpha_i}(X_i)+(\alpha_i-X_{i,1})\bigl(b_i(X)-\alpha_i^{d_i}\bigr).
\]
It vanishes to second order at $p$, its quadratic part is positive definite
on the gate's own coordinates, and its only coupling to earlier gates is the
quadratic cross term $-u_{i,1}\sum_{j,e}a'_{ije}u_{j,e}$ in centered
coordinates $u=X-p$. Geometrically decreasing gate weights make the weighted
sum $\sum_i\rho^{i-1}E_i^*$ positive definite in all directions, including
signed couplings and arbitrary reuse of earlier gates. The coefficients of
$E_i^*$ involve the algebraic number $\alpha_i$, but $E_i^*$ is an explicit
combination, with coefficients polynomial in $\alpha_i$, of rational quadratics
vanishing at $p$. Replacing $\alpha_i$ by a certified rational approximation in
those coefficients therefore keeps the zero exactly while making the gradient
at $p$ as small as required. Lemma~\ref{lem:quartic-realization} then supplies
the quartic and its certificate. No minimal polynomial of a gate value, or of
the field generated by all gate values, is ever expanded; that field has degree
at most $\prod_id_i$.

\begin{remark}[Monotone cube-root circuits]
\label{rem:reductions-monotone}
If $d_i=3$, $c_i\ge1$, and every radicand uses only first powers with
nonnegative coefficients $a_{ij1}\ge0$, boxes need not be part of the input. With
$A=2+\sum_i(c_i+\sum_ja_{ij1})$, the boxes $[1,A]$ pass the direct interval
test, because each radicand interval lies in $[1,A\,(c_i+\sum_ja_{ij1})]
\subseteq[1,A^2]\subseteq[1,A^3]$. Theorem~\ref{thm:reductions-signed-root}
then applies with $\kappa=1$.
\end{remark}

\subsection{Feasibility, a fixed cube, and the sign of the minimum}
\label{sec:reductions-consequences}

Apply Theorem~\ref{thm:reductions-signed-root} to the circuit of
Theorem~\ref{thm:reductions-root-circuit}. All boxes are positive, so every
$\sigma_i=1$ and $\kappa=1/(1-\max_iw_i)\in[1,1+2\cdot10^{-9}]$. The unique zero
has $p_{i,1}=\kappa\xi_i$ and $p_{i,2}=\kappa^2\xi_i^2$, all in $(0,4)$.

\begin{corollary}[Exact feasibility and a fixed cube]
\label{cor:reductions-feasibility}
For the quartic $F$ just described, with designated output gate $o$:
\begin{enumerate}[label=(\roman*)]
\item The set $\{X:\ F(X)\le0,\ X_{o,1}\ge\kappa\}$ is nonempty exactly when
$V>0$, and it is then a single point.
\item Let $X=\Delta Y+\kappa e_{o,1}$, where $\Delta$ is diagonal with entry $8-\kappa$ at
$(o,1)$ and $8$ elsewhere, and put $\widetilde F(Y)=F(\Delta Y+\kappa e_{o,1})$. Then
$\widetilde F$ is a sum of $N+1$ squares of rational quadratics, has a rational
positive definite Hessian Gram, satisfies $\nabla^2\widetilde F\succeq54I$, and
\[
 \min_{Y\in[0,1]^N}\widetilde F(Y)=0\iff V>0;
\]
otherwise the minimum over the cube is positive.
\end{enumerate}
\end{corollary}

\begin{proof}
Since $F\ge0$ vanishes only at $p$, the first set is $\{p\}$ when
$p_{o,1}=\kappa\xi_o\ge\kappa$, that is, when $\xi_o\ge1$, and empty
otherwise. By~\eqref{eq:reductions-compiler-output}, $\xi_o\ge1$ holds exactly
when $\xi_o>1$, that is, when $2V-1>0$. For~(ii), the affine map sends
$[0,1]^N$ onto the box with $X_{o,1}\in[\kappa,8]$ and all other coordinates in
$[0,8]$. All coordinates of $p$ lie in $(0,4)$, so this box contains $p$
exactly when $p_{o,1}\ge\kappa$. If it does, the minimum is $F(p)=0$. If not,
$F$ is positive on the compact box. Composition with an affine map preserves
the rational square factors, and Lemma~\ref{lem:models-gram-curvature}
transforms the Hessian Gram by an invertible rational congruence. Finally
$\nabla^2\widetilde F(Y)=\Delta\,\nabla^2F(X)\,\Delta\succeq\frac32\cdot36\,I$ because
$8-\kappa>6$.
\end{proof}

Part~(i) is a degenerate exact feasibility problem: its feasible set is
empty or a single point, so it has no Slater point. Part~(ii) separates that
degeneracy from the domain: the cube has the interior point
$(\frac12,\ldots,\frac12)$, and only the threshold $\widetilde F\le0$ is
degenerate. In the no case of part~(ii), the positive minimum is not shown to
be at least $2^{-\poly(L)}$.

The unconstrained value test needs a different idea, since an unconstrained
minimum of a sum of squares is never negative. The auxiliary gate $e$ of
Theorem~\ref{thm:reductions-root-circuit} supplies a nonzero signal whose
square is at most one eighth of the output signal's absolute value. A cubic
tilt built from the two signals encodes the output sign in the minimum.

\begin{theorem}[Sign of an unconstrained minimum]
\label{thm:reductions-min-sign}
There is a polynomial-time many-one reduction from $\PosSLP$ to the test
$f^*<0$ for certified quartics, whose output instances $(G,M_G)$ satisfy
$M_G\succeq I$ and $G^*\ne0$. More precisely, $V>0$ implies $G^*<0$, and
$V\le0$ implies $G^*>0$.
\end{theorem}

\begin{proof}
Let $F$, $p$, $\kappa$ and the gates $o\ne e$ be as above, and let $M$ be the
rational positive definite Hessian Gram of $F$, of order $N+N^2$. Put
\[
 u(X)=X_{e,1}-\kappa,\qquad v(X)=X_{o,1}-\kappa,\qquad
 u_0=u(p)=\kappa(\xi_e-1),\qquad v_0=v(p)=\kappa(\xi_o-1).
\]
By~\eqref{eq:reductions-compiler-output} and $1\le\kappa<2$,
\begin{equation}
 u_0\ne0,\qquad v_0\ne0,\qquad \operatorname{sign}v_0=\operatorname{sign}(2V-1),
 \qquad \abs{v_0}\le\tfrac18,\qquad u_0^2\le\tfrac{\kappa}{8}\abs{v_0}\le\tfrac12\abs{v_0}.
 \label{eq:reductions-tilt-signals}
\end{equation}
The cubic $-u^2v$ has
$z^{\mathsf T}\nabla^2(-u^2v)(X)z=2\kappa z_a^2+4\kappa z_az_b-2X_bz_a^2-4X_az_az_b$,
where $a=(e,1)$ and $b=(o,1)$ are distinct indices. On the basis $(z,X\otimes z)$
this biform has the rational symmetric Gram matrix $B$ whose only nonzero
entries are
\[
 B_{z_a,z_a}=2\kappa,\quad B_{z_a,z_b}=B_{z_b,z_a}=2\kappa,\quad
 B_{z_a,X_bz_a}=B_{X_bz_a,z_a}=-1,\quad
 B_{z_b,X_az_a}=B_{X_az_a,z_b}=-2,
\]
and $\norm{B}\le\norm{B}_F=(12\kappa^2+10)^{1/2}<8$. Let
$\mu=\det M/(\tr M)^{N+N^2-1}\le\lambda_{\min}(M)$ and
$\lambda=\lceil9/\mu\rceil$. Then
\[
 G=\lambda F-u^2v,\qquad M_G=\lambda M+B\succeq(\lambda\mu-8)I\succeq I,
\]
so $(G,M_G)$ is a certified quartic with $\nabla^2G\succeq I$; the cubic tilt
leaves the quartic part unchanged. All data have polynomial bit length.

Since $F\ge0$ and $F(p)=0$, we have $\nabla F(p)=0$, hence
\[
 G(p)=-u_0^2v_0,\qquad \nabla G(p)=-2u_0v_0e_a-u_0^2e_b,\qquad
 \norm{\nabla G(p)}^2=4u_0^2v_0^2+u_0^4.
\]
If $V>0$, then $v_0>0$ and $G^*\le G(p)<0$. If $V\le0$, write $t=-v_0>0$.
Strong convexity with modulus one gives
$G(X)\ge G(p)+\nabla G(p)^{\mathsf T}(X-p)+\frac12\norm{X-p}^2
\ge G(p)-\frac12\norm{\nabla G(p)}^2$, so
\[
 G^*\ge u_0^2t-2u_0^2t^2-\tfrac12u_0^4
 =u_0^2\bigl(t-2t^2-\tfrac12u_0^2\bigr)
 \ge u_0^2\bigl(t-\tfrac t4-\tfrac t4\bigr)=\tfrac12u_0^2t>0,
\]
using $t\le\frac18$ and $u_0^2\le\frac t2$ from~\eqref{eq:reductions-tilt-signals}.
\end{proof}

The perturbation can move the minimizer to an irrational point; no rationality
statement is made for $G$. The positive value in the no case is at least
$\frac12u_0^2\abs{v_0}$, a quantity of doubly exponentially small size, and no
better lower bound is claimed. The perturbation argument uses only two
distinct affine signals satisfying~\eqref{eq:reductions-tilt-signals} at the
zero of a certified sum of squares; the compiler is what produces such signals
with polynomial-size data.

\subsection{Sign compiler II: rational unit quaternions}
\label{sec:reductions-quaternion}

Write a quaternion as $u=(u_0,\vec u)$ with scalar part $u_0$ and vector part
$\vec u=(u_1,u_2,u_3)$. The product is
$uw=(u_0w_0-\vec u\cdot\vec w,\ u_0\vec w+w_0\vec u+\vec u\times\vec w)$, the
conjugate is $\bar u=(u_0,-\vec u)$, and $\abs{uw}=\abs u\abs w$ for the
Euclidean norm. For a unit quaternion, $u^{-1}=\bar u$. Let $\mathbf1=(1,0,0,0)$
and $\mathbf i=(0,1,0,0)$.

\begin{definition}[Rational unit-quaternion circuits]
\label{def:reductions-quaternion-circuit}
A \emph{rational unit-quaternion circuit} with $k$ gates lists, in
topological order, gates of three kinds: an explicitly given rational unit
quaternion; the product $g_ag_b$ of two earlier gates, where $a=b$ is allowed;
or the conjugate $\bar g_a$ of an earlier gate. Every gate value is a rational
unit quaternion. The input lists the generators and the operations, not the
expanded coordinates of gate values.
\end{definition}

\begin{theorem}[Quaternion sign compiler]
\label{thm:reductions-quaternion}
Let $\mathbf c=\frac12(1,1,1,1)$ and
$\mathbf w_\vartheta=\bigl(\frac{1-\vartheta^2}{1+\vartheta^2},\frac{2\vartheta}{1+\vartheta^2},0,0\bigr)$
with $\vartheta=2^{-20}$. Given an integer straight-line program with output
$V$, one can compute in polynomial time a rational unit-quaternion circuit
whose generators are among $\mathbf1,\mathbf i,\mathbf c,\mathbf w_\vartheta$,
together with an output gate whose value $u$ satisfies
\[
 u_1\ne0,\qquad \operatorname{sign}u_1=\operatorname{sign}(2V-1).
\]
Consequently, deciding whether a specified coordinate of a specified gate of a
rational unit-quaternion circuit is positive is $\PosSLP$-complete.
\end{theorem}

The proof is in Appendix~\ref{app:reductions-quaternion-compiler}. Two
conjugations do the work. Conjugation by $\mathbf i$ is
$\iota(u)=(u_0,u_1,-u_2,-u_3)$, and conjugation by $\mathbf c$ permutes the axes
cyclically, $\operatorname{rot}(u)=(u_0,u_3,u_1,u_2)$. For a unit quaternion,
the product
\[
 \operatorname{pr}(u)=u\,\iota(u)=\bigl(1-2u_1^2,\ 2u_0u_1,\ 2u_1u_3,\ -2u_1u_2\bigr)
\]
keeps the selected coordinate, doubled, and suppresses the other vector
coordinates to second order. The commutator satisfies the exact identity
$[u,w]-\mathbf1=uw\bar u\bar w-\mathbf1=2(0,\vec u\times\vec w)\,\bar u\bar w$, so
near the identity it multiplies small rotation vectors. Accordingly,
$\operatorname{add}(u,w)=\operatorname{pr}(uw)$ adds two signals and
$\operatorname{prod}(u,w)=\operatorname{pr}(\operatorname{rot}([u,\operatorname{rot}(w)]))$
multiplies them, up to the constant factors two and four and errors of
relative order $\delta$. Iterating $\operatorname{prod}(g,g)$ from
$\mathbf w_\vartheta$ produces an exact rational quaternion whose selected
coordinate $\delta$ satisfies $0<\delta<2^{-16\cdot2^j}$ after $j$ steps. The same numerator--denominator homogenization as in the cube-root
compiler handles orders and cancellation; here the common leading factors are
four for products and eight for sums. The upper bound is elementary: four
integer numerator circuits and one positive common denominator per gate
represent every gate value exactly.

\begin{proposition}[Realization of quaternion circuits]
\label{prop:reductions-quaternion-realization}
Given a rational unit-quaternion circuit with $k$ gates and values
$p_1,\ldots,p_k$, one can compute in polynomial time a rational quartic $F$ in
the $N=4k$ variables $X=(X_1,\ldots,X_k)$, $X_i\in\R^4$, such that $F$ is
supplied as a sum of $N+1$ squares of rational quadratics,
$\nabla^2F\succeq\frac32I$, $F^{-1}(0)=\{p\}$ with $p=(p_1,\ldots,p_k)$, and a
rational positive definite Hessian Gram for $F$ is part of the output. The
coordinates of $p$ are preserved exactly; in particular $p\in\Q^N\cap[-1,1]^N$.
\end{proposition}

The proof is in Appendix~\ref{app:reductions-quaternion-realization}. The
residuals are $X_i-p_i$ for a constant gate, $X_i-X_aX_b$ and $X_i-\bar X_a$. With
$q_i=\abs{X_i}^2-1$, the exposing quadratic of a product gate is
$q_i-2\langle p_i,X_i-X_aX_b\rangle-q_a-q_b$, which equals
$\abs{X_i-p_i}^2-\abs{X_a-p_i\bar X_b}^2$ because $p_i\bar p_b=p_a$ for unit
quaternions; the identity also holds when $a=b$. Weights $16^{1-i}$ absorb the
negative terms of all later gates. Only the coefficients $p_i$ multiplying the
residuals are rounded, so the zero stays exactly at the rational gate list.

\begin{proof}[Proof of Theorem~\ref{thm:rational-optimizer}]
Apply Theorem~\ref{thm:reductions-quaternion} and then
Proposition~\ref{prop:reductions-quaternion-realization}, and let $j$ index the
first vector coordinate of the output gate. The unique zero of $F$ is its
unique minimizer, so $F^*=0$ and $p$ is the rational list of gate values,
which lies in $[-1,1]^N$ because every gate is a unit quaternion. The
selected coordinate satisfies $p_j\ne0$ and $p_j>0$ exactly when $V>0$. The
remaining promises are part of
Proposition~\ref{prop:reductions-quaternion-realization}. This proves the
reduction to $p_j>0$ and, since $p_j\ne0$, also to $p_j\ge0$. Replacing $V$ by
$1-V$ gives the tests $p_j<0$ and $p_j\le0$. The upper bounds follow from
Theorem~\ref{thm:exact-upper} with the observable $x_j$.
\end{proof}

\begin{remark}[Length of the rational optimizer]
\label{rem:reductions-rational-length}
In Theorem~\ref{thm:rational-optimizer}, the generator gate after $2T+2$
iterations is one of the gates, so its selected coordinate is a coordinate of
$p$. It is a nonzero rational of absolute value below $2^{-64\cdot4^T}$
(Lemma~\ref{lem:reductions-generator}), so its reduced denominator has more
than $64\cdot4^T$ bits. The construction has $N=O(T)$ variables, so this
length is exponential in $N$. The reduction never computes this fraction.
\end{remark}

\subsection{Proof of the classification}
\label{sec:reductions-proof-main}

\begin{proof}[Proof of Theorem~\ref{thm:quartic-complete}]
\emph{Validity.} A reduction first checks the format
\eqref{eq:reductions-certified} by coefficient comparison and rational
elimination. On an invalid input it outputs a fixed program with output
$-1$, which is a no-instance of $\PosSLP$. All further statements concern
valid inputs.

\emph{Upper bounds in (i) and (ii).} A valid certified quartic satisfies the
hypotheses of Theorem~\ref{thm:exact-upper}, with the checked Hessian Gram as
curvature certificate. Applying it to the observable $h=f-r$, for which
$h(p)=f^*-r$, handles the five relations between $f^*$ and $r$, and $h=x_j-r$
handles those between $p_j$ and $r$; each yields one $\PosSLP$ instance.

\emph{Lower bounds in (i).} Theorem~\ref{thm:reductions-min-sign} maps an
instance with output $V$ to a certified quartic with $r=0$, $M_G\succeq I$,
$G^*\ne0$, and $V>0\iff G^*<0$. Because $G^*\ne0$, the same instance also
decides $G^*\le0$. Applied to the program with output $1-V$, which is positive
exactly when $V\le0$, it gives $V>0\iff G^*>0\iff G^*\ge0$. The coordinate tests
are Theorem~\ref{thm:rational-optimizer}.

\emph{Part (iii).} For a valid certified quartic, $f-f^*$ is a sum of squares
of real polynomials by Taylor integration of the Hessian Gram at the
stationary point $p$ (Lemma~\ref{lem:taylor-sos}). Hence $f\ge0$, $f^*\ge0$ and
membership in the real sum-of-squares cone are equivalent. If $f$ has a
positive definite Gram matrix $Q$ on the monomial vector $m(x)$ of degree at
most two, then $f(x)\ge\lambda_{\min}(Q)\norm{m(x)}^2\ge\lambda_{\min}(Q)>0$,
because $m(x)$ contains the constant monomial; hence $f^*>0$. Conversely, if
$f^*>0$, a rational positive definite Gram exists by
Theorem~\ref{thm:interior-gram}. Thus these three properties are the tests
$f^*\ge0$, $f^*\ge0$ and $f^*>0$, which reduce to $\PosSLP$ by the upper bounds
in (i). For the lower bounds, apply Theorem~\ref{thm:reductions-min-sign} to the
program with output $1-V$. If $V>0$, its output quartic has $G^*>0$; it is
nonnegative, real SOS, has a rational positive definite polynomial Gram, and is
a sum of squares of rational polynomials, because rational elimination of a
positive definite rational Gram gives positive rational weights and every
positive rational is a sum of rational squares
(Lemma~\ref{lem:models-rational-squares}). If $V\le0$, then $G^*<0$, and $G$ has
none of these properties.

\emph{Part (iv).} The instances used above have exactly the stated additional
properties.
\end{proof}

\subsection{Related comparisons}
\label{sec:reductions-comparisons}

\paragraph{Sums of numbers with one real conjugate.}
A different source problem is also encoded by one quartic and one affine
inequality. It concerns algebraic numbers given by dense minimal polynomials,
rather than circuits, and its relation to $\PosSLP$ and to the Square Root Sum
problem is not settled here.

\begin{proposition}[One-real-conjugate sums]
\label{prop:reductions-sums}
Let $P_1,\ldots,P_m\in\Q[T]$ be given by dense coefficient lists, each
irreducible with exactly one real root $\alpha_i$, and let
$c_1,\ldots,c_m,b\in\Q$. In polynomial time one can compute a rational quartic
$F$, supplied as a sum of squares of rational quadratics with
$\nabla^2F\succeq I$, and a rational affine function $\ell$, such that
$\{F\le0\}$ is a single point and
\[
 \{F\le0,\ \ell\le b\}\ne\emptyset\iff\sum_{i=1}^mc_i\alpha_i\le b.
\]
The comparison $\sum_ic_i\alpha_i\le b$, and likewise $<$ and $=$, reduces to
$\PosSLP$.
\end{proposition}

\begin{proof}
For each $i$ with $\deg P_i>1$, Theorem~\ref{thm:singleton-field} constructs,
in time polynomial in the dense encoding of $P_i$, a rational quartic $F_i$ on
its own block of variables, given as a sum of squares of rational quadratics,
with $\nabla^2F_i\succeq I$ and a unique zero whose first coordinate is
$\alpha_i$. For a rational root use the one-variable quartic
$(z-\alpha_i)^2+(z-\alpha_i)^4$. Put $F=\sum_iF_i$ on disjoint blocks and
$\ell=\sum_ic_iz_{i,1}$, where $z_{i,1}$ is the first coordinate of block $i$.
Every summand is nonnegative, so $F\le0$ forces every block to its zero, and
the affine row then tests the sum. The Hessian is block diagonal with blocks at
least $I$. For an empty list, use $F(z)=z^2+z^4$ and $\ell=0$. Finally, $F$ has
the supplied curvature bound $\nabla^2F\succeq I$, and the observable $\ell-b$
at its unique minimizer is $\sum_ic_i\alpha_i-b$; Theorem~\ref{thm:exact-upper}
reduces each relation of that observable to zero to one $\PosSLP$
instance.
\end{proof}

Signed sums of real cube roots of integers are the special case
$P_i=T^3-a_i$. Such a cubic is irreducible unless $a_i$ is a perfect cube,
which binary search detects in polynomial time; negative radicands only change
the sign of $c_i$. Unlike Theorem~\ref{thm:quartic-complete}, the proposition is
an upper bound for its source problem together with a convex encoding of it.
We do not know whether $\PosSLP$ reduces to the comparison of sums of
one-real-conjugate numbers, and none of our results implies such a reduction.
Square Root Sum is not covered by this construction: every coordinate of a
singleton defined by rational convex polynomial inequalities has exactly one
real conjugate (Theorem~\ref{thm:singleton-field}), so no such singleton has a
coordinate $\sqrt2$. This obstruction concerns value-preserving encodings only;
it does not exclude other reductions. The block sum has a rational positive
semidefinite Hessian Gram, obtained by embedding the block certificates. When
$m\ge2$, so that $F$ has at least two nonempty blocks of variables, it has no
positive definite Hessian Gram on the full basis
(Remark~\ref{rem:models-gram-kinds}). With a single block, and for the empty
list, $F$ is itself a quartic of Theorem~\ref{thm:singleton-field}, so it has
a positive definite rational Hessian Gram; for the empty list, $z^2+z^4$ has
the Hessian Gram $\diag(2,12)$ on $(v,zv)$. The reduction to $\PosSLP$ in
the proof uses the supplied curvature bound $\nabla^2F\succeq I$, which holds
in every case, and needs no positive definite Hessian Gram.

\paragraph{A compact convex program with the same optimizer.}
Placing the values of a certified root circuit at the unique optimizer of
\emph{some} compact convex program is much easier than
Theorem~\ref{thm:reductions-signed-root}.

\begin{proposition}[Box-convex baseline]
\label{prop:reductions-box-baseline}
Let a certified odd-root circuit have boxes with $\underline\xi_i\ge1$, which
the normalization in Theorem~\ref{thm:reductions-signed-root} always achieves,
and let $A\ge1$ bound every $\overline\xi_i$ and the sum of the absolute values
of all coefficients. In the variables $X_i$ and $Y_{i,e}$, $2\le e\le n_i$, with
$Y_{i,1}=X_i$, consider the compact convex set defined by the boxes
$X_i\in[\underline\xi_i,\overline\xi_i]$,
$Y_{i,e}\in[\underline\xi_i^{\,e},\overline\xi_i^{\,e}]$ and the inequalities
\[
 X_i^{d_i}\le c_i+\sum_{j<i,\,e}a_{ije}Y_{j,e},\qquad X_i^e\le Y_{i,e}
 \quad(2\le e\le n_i).
\]
Order the variables gate by gate, each root before its powers, and set
$y_t=X_i$ or $y_t=-Y_{i,e}$ accordingly. With $C=NA^N$ and
$W_t=(2C+2)^{N-t}$, the circuit point $X_i=\xi_i$, $Y_{i,e}=\xi_i^e$ is the
unique maximizer of $\sum_tW_ty_t$ over this set. All data have polynomial bit
length.
\end{proposition}

The proof is in Appendix~\ref{app:reductions-comparisons}. The program has
constraints of the circuit's own degrees, needs the boxes to be convex, and
represents the circuit only as the optimizer of a linear objective. The
realization of Theorem~\ref{thm:reductions-signed-root} adds what the
comparisons in Theorem~\ref{thm:quartic-complete} need: a single globally
strongly convex polynomial of degree four, without constraints, with the
circuit as its unique zero and a checked rational Hessian certificate.

\subsection{Scope}
\label{sec:reductions-scope}

The lower bounds are conditional statements about exact computation: a
polynomial-time algorithm for any problem in Theorem~\ref{thm:quartic-complete}(i)
or (iii) would give one for $\PosSLP$. They imply neither NP-hardness nor a
superpolynomial lower bound, since neither is known for $\PosSLP$, and they give
no lower bound for approximate value or point computation. Equality tests have only the upper bound of
Theorem~\ref{thm:exact-upper}; the constructed instances deliberately avoid
equality. Every hard instance encodes its sign by a doubly exponentially small
signal, so no inverse-polynomial separation of the hard instances from the
threshold is asserted. In particular, the results say nothing about
algorithms that may err when the optimal value or coordinate lies within
$2^{-\poly(L)}$ of the threshold. The quartics are explicitly expanded with polynomially
many monomials; the lower bounds do not depend on succinct or circuit-encoded
input. They concern unconstrained optimization, a halfspace, or a fixed cube.
Exact comparison over general rational polyhedra is treated in
Section~\ref{sec:constraints}.
